\section{Rasprava}

Rezultati pokazuju da ispitanici najčešće prepoznaju sigurnosne značajke kao najveću prednost SMR tehnologije. Taj odgovor odabralo je 33,86\% ispitanika, dok 29,13\% ispitanika kao glavnu prednost navodi kraće vrijeme izgradnje i puštanja u pogon. Niže početne i ukupne troškove investicije odabralo je 16,19\% ispitanika, a 20,82\% ispitanika odgovorilo je da ne zna. Takva raspodjela upućuje na to da se SMR tehnologija u percepciji ispitanika ne promatra isključivo kroz ekonomske argumente, nego znatan dio ispitanika prednost vidi upravo u sigurnosnim i izvedbenim obilježjima tehnologije.

Kod pitanja o državnom subvencioniranju razvoja SMR tehnologije vidljiva je većinska potpora: 59,01\% ispitanika smatra subvenciju opravdanom, 8,01\% smatra da nije opravdana, dok 32,98\% ispitanika nije sigurno. Visok udio odgovora \textit{Ne znam} važan je nalaz jer pokazuje da, uz većinsku načelnu potporu, postoji i znatna razina nesigurnosti ili nedovoljne formiranosti stavova o ekonomskim i energetskim posljedicama takve javne politike.

Demografske usporedbe pokazale su statistički značajnu povezanost između spola i odgovora na pitanje P9 (p < 0,001, Cramérov \(V = 0{,}283\)). Prema grafičkom prikazu, muškarci češće ističu kraću izgradnju kao najveću prednost SMR tehnologije, dok žene češće ističu sigurnosne značajke ili odgovaraju da ne znaju. Ova povezanost je sadržajno relevantna jer upućuje na različite naglaske u percepciji tehnologije: dio ispitanika SMR promatra kroz izvedbene i vremenske prednosti, dok drugi dio veću važnost pridaje sigurnosti ili pokazuje veću nesigurnost.

Povezanost između obrazovanja i odgovora na pitanje P10 također je statistički značajna (p = 0,01428), ali je Cramérov \(V = 0{,}100\), što upućuje na slabu povezanost. Dodatno, kod tog testa dio očekivanih frekvencija ostaje manji od 5, pa rezultat treba interpretirati oprezno. Iako se može uočiti da se udjeli potpore subvencijama razlikuju među obrazovnim skupinama, ti se nalazi ne bi trebali prenaglašavati bez dodatnog uzorka ili bolje uravnoteženih kategorija obrazovanja.

Najizraženija povezanost u analizi javlja se između odgovora na P9 i P10 (p < 0,001, Cramérov \(V = 0{,}318\)). Ispitanici koji kao najveću prednost SMR tehnologije ističu sigurnosne značajke najčešće podržavaju državne subvencije, dok ispitanici koji na P9 odgovaraju \textit{Ne znam} u najvećem udjelu također nisu sigurni ni oko opravdanosti subvencija. To sugerira da percepcija konkretne prednosti tehnologije može biti povezana s većom spremnošću na podršku javnom ulaganju u njezin razvoj. Drugim riječima, ispitanici koji imaju jasniju pozitivnu predodžbu o SMR tehnologiji češće podržavaju i aktivniju državnu ulogu u razvoju te tehnologije.

Pri interpretaciji treba uzeti u obzir da anketa mjeri stavove ispitanika, a ne tehničku ili ekonomsku opravdanost pojedinih energetskih rješenja. Također, statistička značajnost ne mora nužno značiti i veliku praktičnu važnost, zbog čega se uz p-vrijednosti promatra i Cramérov \(V\).

Ograničenje analize predstavlja i struktura uzorka. Dobna skupina od 15 do 24 godine čini najveći dio ispitanika, pa rezultati ne moraju u potpunosti predstavljati stavove ukupne populacije. Osim toga, odgovori na pitanje o županiji nisu standardizirani, zbog čega se kod tog obilježja pojavljuje velik broj rijetkih kategorija i varijanti zapisa. Zbog toga su rijetke kategorije prije testiranja spojene u skupinu \textit{Ostalo (<5)}, ali rezultate za županiju ipak treba smatrati manje pouzdanima od rezultata za jasnije definirana obilježja poput spola i dobi.
